\chapter{拓扑空间}
\label{chap:topological-spaces}

度量把两点之间的接近程度表示为一个非负实数，但收敛与连续的许多论证并不使用这个数值本身。它们只使用哪些集合能够作为点的邻域，以及原像是否保持这类集合。函数空间的弱拓扑、任意多个空间的乘积拓扑和商空间都说明，仅靠一个预先给定的距离来组织开集并不总是自然。

拓扑空间保留开集及其运算规则，不要求开集来自距离。本章建立这种语言，并特别说明两件在度量空间中并不突出的事实：拓扑可以由一族映射反向规定；在一般拓扑空间中，数列不一定能够检测闭包与连续性。

\section{开集公理}

\begin{definition}[拓扑与拓扑空间]
设 \(X\) 为集合。若子集族 \(\mathcal T\subseteq\mathcal P(X)\) 满足
\begin{enumerate}[label=\textup{(\roman*)}]
 \item \(\varnothing,X\in\mathcal T\)；
 \item \(\mathcal T\) 中任意多个集合的并仍属于 \(\mathcal T\)；
 \item \(\mathcal T\) 中有限多个集合的交仍属于 \(\mathcal T\)，
\end{enumerate}
则称 \(\mathcal T\) 是 \(X\) 上的一个\term{拓扑}，称二元组
\((X,\mathcal T)\) 为\term{拓扑空间}。\(\mathcal T\) 的成员称为开集。
\end{definition}

“任意并、有限交”的不对称来自距离空间中的开球。任意多个开集的并在每一点仍能找到一个来自其中某个开集的小球；无限多个开集的交却可能把可用半径压到零，例如
\[
 \bigcap_{n=1}^{\infty}(-1/n,1/n)=\{0\}.
\]

同一集合可以带有不同拓扑。\(\{\varnothing,X\}\) 称为不可分拓扑，\(\mathcal P(X)\) 称为离散拓扑。前者只记录“空”与“全体”，后者把每个子集都视为开集。若 \(d\) 是 \(X\) 上的度量，由开球生成的开集族给出度量拓扑。拓扑语言因而包含前一编的全部度量空间。

\begin{example}[余有限拓扑]
设 \(X\) 为无限集，令
\[
 \mathcal T_{\mathrm{cof}}
 =\{\varnothing\}\cup\{U\subseteq X:X\setminus U\text{ 有限}\}.
\]
任意并的补集是有限集族的交，仍为有限集；有限交的补集是有限并，也为有限集。因此它是一个拓扑。它在 \(X\) 不可数时不能由通常那样的“小半径”直观刻画，却是检验分离性、紧致性和数列论证的标准例子。
\end{example}

若 \(\mathcal T_1\subseteq\mathcal T_2\)，称 \(\mathcal T_1\) 比 \(\mathcal T_2\) 粗，或 \(\mathcal T_2\) 比 \(\mathcal T_1\) 细。拓扑越细，可供检验连续性的开集越多；同一个由 \(X\) 出发的映射通常越难连续。

\section{闭集与邻域}

\begin{definition}[闭集与邻域]
拓扑空间 \(X\) 的子集 \(F\) 称为闭集，如果 \(X\setminus F\) 开。若 \(x\in X\)，集合 \(N\subseteq X\) 称为 \(x\) 的邻域，如果存在开集 \(U\) 使
\[
 x\in U\subseteq N.
\]
邻域本身不要求开。
\end{definition}

由 De Morgan 律，闭集族包含 \(\varnothing,X\)，在任意交与有限并下封闭。反过来，任何满足这三条规则的子集族取补以后都给出一个拓扑。

前一编定义的内部、闭包和边界不依赖距离。对 \(A\subseteq X\)，令
\[
 \operatorname{int}A=\bigcup\{U\in\mathcal T:U\subseteq A\},
 \qquad
 \overline A=\bigcap\{F:F\text{ 闭且 }A\subseteq F\},
\]
并令 \(\partial A=\overline A\cap\overline{X\setminus A}\)。于是 \(x\in\overline A\) 当且仅当 \(x\) 的每个邻域都与 \(A\) 相交。这个邻域判据以后比“取趋于 \(x\) 的数列”更可靠。

\begin{proposition}[邻域系统刻画]
\label{prop:neighborhood-system-topology}
对每个 \(x\in X\)，记其邻域族为 \(\mathcal N(x)\)。则
\begin{enumerate}[label=\textup{(\roman*)}]
 \item 每个 \(N\in\mathcal N(x)\) 都含 \(x\)；
 \item \(N\in\mathcal N(x)\)、\(N\subseteq M\) 蕴含 \(M\in\mathcal N(x)\)；
 \item 两个 \(x\) 的邻域之交仍是 \(x\) 的邻域；
 \item 对每个 \(N\in\mathcal N(x)\)，存在 \(M\in\mathcal N(x)\) 使对所有 \(y\in M\)，都有 \(N\in\mathcal N(y)\)。
\end{enumerate}
反之，若一族 \(\mathcal N(x)\) 满足这四条，则
\[
 U\text{ 开}\quad\Longleftrightarrow\quad
 U\in\mathcal N(x)\text{ 对每个 }x\in U
\]
定义了 \(X\) 上唯一的拓扑，并以原来的 \(\mathcal N(x)\) 为邻域族。
\end{proposition}

\begin{proof}
由开集定义，前四条逐项成立；第四条可取 \(x\in V\subseteq N\) 的开集 \(V\)，再令 \(M=V\)。

反过来，把满足右侧条件的集合称为开集。空集与全体满足要求，任意并也满足。若 \(U,V\) 开且 \(x\in U\cap V\)，则 \(U,V\in\mathcal N(x)\)，故 \(U\cap V\in\mathcal N(x)\)，所以有限交开。

还需核对给定的邻域族没有改变。若 \(N\in\mathcal N(x)\)，由(iv)以及必要时把所得邻域与 \(N\) 相交，可取 \(M\in\mathcal N(x)\)、\(M\subseteq N\)，且 \(N\in\mathcal N(y)\) 对每个 \(y\in M\) 成立。令
\[
 U=\{y\in N:N\in\mathcal N(y)\}.
\]
则 \(x\in U\subseteq N\)。对 \(y\in U\)，再用(iv)可在 \(U\) 内找到 \(y\) 的邻域，故 \(U\) 开，于是 \(N\) 是新拓扑下 \(x\) 的邻域。反向包含由(i)、(ii)立即得到。
\end{proof}

\section{基和子基}

逐个列出全部开集往往既不经济，也不能显示拓扑从何而来。度量空间只需列出开球；乘积空间只需列出有限多个坐标条件。基和子基把这种做法抽象出来。

\begin{definition}[拓扑基]
设 \(\mathcal B\subseteq\mathcal P(X)\)。若
\begin{enumerate}[label=\textup{(\roman*)}]
 \item 对每个 \(x\in X\)，存在 \(B\in\mathcal B\) 使 \(x\in B\)；
 \item 若 \(x\in B_1\cap B_2\)，其中 \(B_1,B_2\in\mathcal B\)，则存在 \(B_3\in\mathcal B\) 使
 \(x\in B_3\subseteq B_1\cap B_2\)，
\end{enumerate}
则称 \(\mathcal B\) 为 \(X\) 上的一个\term{拓扑基}。由 \(\mathcal B\) 生成的开集是所有基元素的并。
\end{definition}

第二条不是要求 \(B_1\cap B_2\) 本身属于基，而是要求交集内每一点都有更小的基元素。它恰好保证基元素之并在有限交下封闭。

\begin{proposition}[判定一族集合是否为给定拓扑的基]
\label{prop:basis-neighborhood-criterion}
设 \(X\) 已带拓扑 \(\mathcal T\)，且 \(\mathcal B\subseteq\mathcal T\)。则 \(\mathcal B\) 是 \(\mathcal T\) 的基，当且仅当对每个开集 \(U\) 和每个 \(x\in U\)，存在 \(B\in\mathcal B\) 使
\[
 x\in B\subseteq U.
\]
\end{proposition}

\begin{proof}
若 \(\mathcal B\) 是基，把 \(U\) 写成基元素之并，含 \(x\) 的某一项即满足要求。反之，对每个 \(x\in U\) 选取相应 \(B_x\)，便有 \(U=\bigcup_{x\in U}B_x\)。所以每个开集都由基生成；而基元素原已开，两种拓扑相同。
\end{proof}

\begin{definition}[子基]
覆盖 \(X\) 的子集族 \(\mathcal S\) 称为一个\term{子基}。它的有限交（包括空交 \(X\)）组成一个基；该基生成的拓扑称为由 \(\mathcal S\) 生成的拓扑。
\end{definition}

例如，实直线通常拓扑可由开区间组成的基生成，也可由射线
\((-\infty,a)\)、\((a,\infty)\) 组成的子基生成。子基在
\chapterref{chap:compact-topological-spaces}的乘积紧致性中还有更深的作用：只检验子基本闭集的有限交性质便足以推出紧致性。

\section{可数性公理}

\begin{definition}[第一、第二可数性]
若每个 \(x\in X\) 都有至多可数的邻域基，即存在邻域列 \((N_n)\)，使每个 \(x\) 的邻域都包含某个 \(N_n\)，则称 \(X\) 满足第一可数性公理。若 \(X\) 的拓扑有一个至多可数的基，则称 \(X\) 满足第二可数性公理。
\end{definition}

度量空间在每点以 \(B(x,1/n)\) 为邻域基，因而第一可数。可分度量空间还是第二可数：若 \(D\) 为至多可数稠密集，则以 \(D\) 为中心、正有理数为半径的开球组成可数基。一般拓扑空间中，可分性、第一可数性和第二可数性不再等价。

\begin{proposition}[第二可数性的两个后果]
\label{prop:second-countable-consequences}
第二可数空间是第一可数且可分的；并且它的每个开覆盖都有至多可数的子覆盖。
\end{proposition}

\begin{proof}
设 \(\mathcal B=\{B_n:n\in\N\}\) 是可数基。对固定 \(x\)，所有含 \(x\) 的 \(B_n\) 构成可数邻域基。对每个非空 \(B_n\) 选一点 \(x_n\in B_n\)，所得至多可数集 \(D\) 与每个非空开集相交，故稠密。

若 \(\mathcal U\) 为开覆盖，对每个被某个 \(U\in\mathcal U\) 包含的基元素 \(B_n\)，选取一个这样的 \(U_n\)。任给 \(x\in X\)，先取 \(U\in\mathcal U\) 含 \(x\)，再取 \(x\in B_n\subseteq U\)，于是 \(x\in U_n\)。故 \(\{U_n\}\) 是至多可数子覆盖。
\end{proof}

这里两次选择都只在至多可数族中进行。若希望完全追踪选择原则，可以固定基的枚举，并在每一步取满足条件的最小指标；从任意开覆盖中选取成员仍使用可数选择的一个形式。本书工作在 ZFC 背景下。

第一可数性恢复了熟悉的序列判据。

\begin{theorem}[第一可数空间中的闭包]
\label{thm:first-countable-sequential-closure}
设 \(X\) 第一可数，\(A\subseteq X\)。则
\[
 x\in\overline A
 \quad\Longleftrightarrow\quad
 \text{存在 }(a_n)\subseteq A\text{ 使 }a_n\to x.
\]
\end{theorem}

\begin{proof}
反向由邻域定义立即成立。对正向，取 \(x\) 的可数邻域基 \((N_n)\)，并用有限交替换为递减邻域基
\(M_n=N_0\cap\cdots\cap N_n\)。因 \(x\in\overline A\)，可取 \(a_n\in A\cap M_n\)。若 \(N\) 是 \(x\) 的邻域，某个 \(M_k\subseteq N\)；当 \(n\ge k\) 时 \(a_n\in M_n\subseteq M_k\subseteq N\)，故 \(a_n\to x\)。
\end{proof}

\section{连续映射与同胚}

\begin{definition}[连续映射]
设 \(X,Y\) 为拓扑空间。映射 \(f:X\to Y\) 称为连续，如果 \(Y\) 的每个开集 \(V\) 的原像 \(f^{-1}(V)\) 都在 \(X\) 中开。
\end{definition}

原像与任意并、有限交及补集相容，因此开集定义同时给出闭集版本。也可用邻域表述：\(f\) 在 \(x\) 连续，当且仅当 \(f(x)\) 的每个邻域 \(V\) 都有 \(x\) 的邻域 \(U\) 满足 \(f(U)\subseteq V\)。

\begin{proposition}[连续性的若干等价形式]
\label{prop:topological-continuity-equivalences}
对映射 \(f:X\to Y\)，下列命题等价：
\begin{enumerate}[label=\textup{(\roman*)}]
 \item \(f\) 连续；
 \item \(Y\) 中每个闭集的原像闭；
 \item 对每个 \(A\subseteq X\)，有 \(f(\overline A)\subseteq\overline{f(A)}\)；
 \item 对 \(Y\) 的某个基 \(\mathcal B\)，每个 \(B\in\mathcal B\) 的原像开。
\end{enumerate}
\end{proposition}

\begin{proof}
(i)与(ii)由取补等价。(i)推出(iii)：若 \(x\in\overline A\) 且 \(V\) 是 \(f(x)\) 的邻域，则 \(f^{-1}(V)\) 是 \(x\) 的邻域，故与 \(A\) 相交，于是 \(V\cap f(A)\ne\varnothing\)。

(iii)推出(ii)：若 \(F\subseteq Y\) 闭，取 \(A=f^{-1}(F)\)，则
\[
 f(\overline A)\subseteq\overline{f(A)}\subseteq F,
\]
故 \(\overline A\subseteq f^{-1}(F)=A\)，所以 \(A\) 闭。(i)与(iv)则由每个开集都是基元素的并得到。
\end{proof}

双射 \(f:X\to Y\) 若 \(f\) 与 \(f^{-1}\) 都连续，称为同胚。同胚不保留距离或角度，而保留开集、闭包、紧致性与连通性等由拓扑定义的性质。

\section{初始拓扑与终拓扑}

\subsection{由出发映射规定拓扑}

设 \(X\) 只是一个集合，已有一族映射 \(f_i:X\to Y_i\)，其中各 \(Y_i\) 已带拓扑。若希望这些映射连续，就必须把每个 \(f_i^{-1}(V)\) 规定为开集。除此之外不加入额外开集，便得到最粗的可用拓扑。

\begin{definition}[初始拓扑]
映射族 \((f_i:X\to Y_i)_{i\in I}\) 在 \(X\) 上诱导的\term{初始拓扑}，是由
\[
 \{f_i^{-1}(V):i\in I,\ V\subseteq Y_i\text{ 开}\}
\]
作为子基生成的拓扑。
\end{definition}

\begin{theorem}[初始拓扑的泛性质]
\label{thm:initial-topology-universal-property}
给 \(X\) 配备上述初始拓扑。对任意拓扑空间 \(Z\) 和映射 \(g:Z\to X\)，
\[
 g\text{ 连续}
 \quad\Longleftrightarrow\quad
 f_i\circ g\text{ 对每个 }i\in I\text{ 连续}.
\]
\end{theorem}

\begin{proof}
正向由连续映射复合保持连续。反向只需检验子基原像，而
\[
 g^{-1}(f_i^{-1}(V))=(f_i\circ g)^{-1}(V)
\]
开。
\end{proof}

\subsection{由到达映射规定拓扑}

方向反过来，给定映射族 \(q_i:X_i\to Y\)，希望它们都连续时，\(Y\) 上能取的最细拓扑也有直接描述。

\begin{definition}[终拓扑]
由 \((q_i:X_i\to Y)_{i\in I}\) 在 \(Y\) 上诱导的\term{终拓扑}定义为
\[
 U\subseteq Y\text{ 开}
 \quad\Longleftrightarrow\quad
 q_i^{-1}(U)\text{ 在 }X_i\text{ 中开，对每个 }i\in I.
\]
\end{definition}

\begin{theorem}[终拓扑的泛性质]
\label{thm:final-topology-universal-property}
给 \(Y\) 配备终拓扑。对任意拓扑空间 \(Z\) 和映射 \(h:Y\to Z\)，
\[
 h\text{ 连续}
 \quad\Longleftrightarrow\quad
 h\circ q_i\text{ 对每个 }i\in I\text{ 连续}.
\]
\end{theorem}

\begin{proof}
正向仍由复合得到。反向取 \(Z\) 中开集 \(V\)。对每个 \(i\)，
\[
 q_i^{-1}(h^{-1}(V))=(h\circ q_i)^{-1}(V)
\]
开，故按终拓扑定义 \(h^{-1}(V)\) 开。
\end{proof}

\section{子空间、积空间和商空间}

\subsection{子空间与积空间}

若 \(A\subseteq X\)，包含映射 \(\iota:A\hookrightarrow X\) 诱导的初始拓扑由
\[
 \{A\cap U:U\subseteq X\text{ 开}\}
\]
组成，称为子空间拓扑。

对空间族 \((X_i)_{i\in I}\)，笛卡儿积 \(\prod_{i\in I}X_i\) 关于坐标投影 \(\pi_i\) 的初始拓扑称为积拓扑。它的一个基由有限坐标柱集
\[
 \prod_{i\in I}U_i,
 \qquad U_i\subseteq X_i\text{ 开，且除有限多个 }i\text{ 外 }U_i=X_i
\]
组成。有限坐标限制十分重要；若允许每个坐标都取任意开集，得到的是通常更细的箱拓扑。

由初始拓扑的泛性质立即得到：映射 \(g:Z\to\prod_iX_i\) 连续，当且仅当每个坐标映射 \(\pi_i\circ g\) 连续。

\subsection{商拓扑与映射下降}

若 \(q:X\to Y\) 是满射，\(Y\) 关于单个映射 \(q\) 的终拓扑称为商拓扑，\(q\) 称为商映射。子集 \(U\subseteq Y\) 开，当且仅当 \(q^{-1}(U)\) 开。终拓扑泛性质化为
\[
 h:Y\to Z\text{ 连续}
 \quad\Longleftrightarrow\quad
 h\circ q:X\to Z\text{ 连续}.
\]

\begin{proposition}[商映射的下降]
\label{prop:quotient-map-descent}
设 \(q:X\to Y\) 为商映射，\(g:X\to Z\) 连续，并且
\[
 q(x)=q(x')\Longrightarrow g(x)=g(x').
\]
则存在唯一连续映射 \(\widetilde g:Y\to Z\) 使 \(g=\widetilde g\circ q\)。
\end{proposition}

\begin{proof}
令 \(\widetilde g(q(x))=g(x)\)。纤维上常值保证良定义，\(q\) 满射保证唯一性。由终拓扑泛性质，\(\widetilde g\) 连续当且仅当 \(\widetilde g\circ q=g\) 连续。
\end{proof}

\section{数列刻画失效的例子}

在度量空间中，闭包、闭集和连续性都可由数列刻画。其证明使用每点的可数邻域基。下面的乘积空间没有这种局部可数性。

\begin{example}[闭包中没有序列极限]
\label{ex:uncountable-product-sequential-failure}
取不可数指标集 \(I\)，令
\[
 X=\{0,1\}^{I}
\]
带积拓扑，其中每个 \(\{0,1\}\) 带离散拓扑。记
\[
 A=\{x\in X:\operatorname{supp}x=\{i:x_i=1\}\text{ 有限}\},
 \qquad p=(1)_{i\in I}.
\]
任意含 \(p\) 的基本开集只限制有限多个坐标。令 \(x\) 在这些坐标上取 \(1\)，其余坐标取 \(0\)，便有 \(x\in A\) 且落在该基本开集内。因此 \(p\in\overline A\)。

然而不存在 \(A\) 中数列收敛到 \(p\)。若 \((x^{(n)})\subseteq A\)，集合
\[
 J=\bigcup_{n\in\N}\operatorname{supp}x^{(n)}
\]
至多可数，故可取 \(i_0\in I\setminus J\)。所有 \(x^{(n)}_{i_0}=0\)，而 \(p_{i_0}=1\)，所以该数列不在坐标 \(i_0\) 上收敛到 \(p\)。
\end{example}

同一论证说明 \(p\) 没有可数邻域基：任给可数族基本邻域，它们涉及的坐标总并仍至多可数；在总并之外增加一个坐标限制，就得到不包含其中任何一个的更小邻域。因此这个乘积空间不是第一可数的。

对任意 \(A\subseteq X\)，把所有可由 \(A\) 中数列收敛得到的极限组成的集合称为 \(A\) 的序列闭包。第一可数空间中它等于 \(\overline A\)；上例中两者不同。

在一般拓扑空间中，“每个收敛数列都满足某性质”只能看到序列可到达的部分。\chapterref{chap:nets-filters}将用网和滤子恢复闭包、紧致性与连续性的无条件收敛刻画。

\section*{习题}
\addcontentsline{toc}{section}{习题}

\subsection*{A组　拓扑、闭包与基}
\begin{enumerate}[label=\textbf{\thechapter.\arabic*.}]
 \exerciseitem{topology-intersection-lattice}
 证明同一集合 \(X\) 上任意一族拓扑的交仍是拓扑。说明两拓扑的并未必是拓扑，并求包含它们的最粗拓扑。
 \exerciseitem{cofinite-closure-convergence}
 设 \(X\) 为无限集并带余有限拓扑。求任意 \(A\subseteq X\) 的闭包，并刻画数列收敛到给定点的条件。
 \exerciseitem{neighborhood-interior-closure}
 只用邻域语言证明
 \[
 x\in\operatorname{int}A\Longleftrightarrow A\in\mathcal N(x),
 \qquad
 x\in\overline A\Longleftrightarrow
 N\cap A\ne\varnothing\quad(N\in\mathcal N(x)).
 \]
 \exerciseitem{basis-comparison}
 设 \(\mathcal B_1,\mathcal B_2\) 分别生成拓扑 \(\mathcal T_1,\mathcal T_2\)。证明
 \(\mathcal T_1\subseteq\mathcal T_2\) 当且仅当每个 \(B\in\mathcal B_1\) 及每个 \(x\in B\) 都有 \(C\in\mathcal B_2\) 满足 \(x\in C\subseteq B\)。
 \exerciseitem{rational-interval-basis}
 证明端点为有理数的开区间构成 \(\R\) 通常拓扑的可数基。
 \exerciseitem{particular-point-countability}
 在不可数集合 \(X\) 中固定 \(p\)，规定非空开集必须含 \(p\)。证明所得空间可分且第一可数，但不第二可数。
\end{enumerate}

\subsection*{B组　连续性与诱导拓扑}
\begin{enumerate}[label=\textbf{\thechapter.\arabic*.},resume]
 \exerciseitem{sierpinski-open-classifier}
 令 \(S=\{0,1\}\) 带拓扑 \(\{\varnothing,\{1\},S\}\)。证明 \(U\subseteq X\) 开，当且仅当其示性函数 \(\mathbf1_U:X\to S\) 连续。
 \exerciseitem{continuity-basis-subbasis}
 证明检验映射 \(f:X\to Y\) 连续时，只需检验 \(Y\) 的一个基，甚至只需检验一个子基的原像为开集。
 \exerciseitem{second-countable-hereditary-product}
 证明第二可数性对子空间遗传，并证明至多可数多个第二可数空间的积仍第二可数。
 \exerciseitem{initial-topology-composition}
 证明初始拓扑构造满足复合律：由 \((g_i\circ f:X\to Z_i)\) 直接诱导的拓扑，等于先在 \(Y\) 上取 \((g_i:Y\to Z_i)\) 的初始拓扑，再由 \(f:X\to Y\) 诱导的拓扑。
 \exerciseitem{product-coordinate-continuity}
 从初始拓扑泛性质证明 \(h:Z\to X\times Y\) 连续当且仅当两个坐标函数连续，并写出有限积的一般版本。
 \exerciseitem{box-versus-product}
 在 \(\R^{\N}\) 中证明 \(\prod_{n\in\Npos}(-1/n,1/n)\) 在箱拓扑中开、在积拓扑中不开。
 \exerciseitem{quotient-saturated-open}
 设 \(q:X\to Y\) 为商映射。称 \(A\subseteq X\) 饱和，如果 \(A=q^{-1}(q(A))\)。证明 \(U\subseteq Y\) 开当且仅当 \(q^{-1}(U)\) 是开饱和集，并说明为什么“\(q\) 把开集映成开集”不是商映射定义的一部分。
 \exerciseitem{circle-as-quotient}
 令 \(q:[0,1]\to S^1\)、\(q(t)=(\cos2\pi t,\sin2\pi t)\)。利用
 \chapterref{chap:compact-topological-spaces}中“紧致空间到 Hausdorff 空间的连续双射为同胚”的结论，证明把 \(0,1\) 识别所得商空间同胚于圆周。
\end{enumerate}

\subsection*{C组　序列方法的范围}
\begin{enumerate}[label=\textbf{\thechapter.\arabic*.},resume]
 \exerciseitem{first-countable-sequential-continuity}
 证明若 \(X\) 第一可数，则 \(f:X\to Y\) 连续当且仅当 \(x_n\to x\) 总蕴含 \(f(x_n)\to f(x)\)。
 \exerciseitem{uncountable-cube-no-countable-base}
 完成正文不可数乘积例的邻域基论证：证明 \(\{0,1\}^{I}\) 在任一点都不第一可数。
 \exerciseitem{finite-support-sequential-closure}
 在正文的 \(A\subseteq\{0,1\}^{I}\) 中，证明 \(A\) 的序列闭包仍为“支集至多可数”的点集；比较它与 \(\overline A=X\)。
\end{enumerate}
